\documentclass[a4paper,11pt]{article}

\usepackage[utf8]{inputenc}
\usepackage[T1]{fontenc}
\usepackage[portuguese]{babel}
\usepackage{amsmath, amssymb}
\usepackage{geometry}
\geometry{margin=2cm}
\usepackage{booktabs}
\usepackage{tabularx}
\usepackage{enumitem}
\usepackage{parskip} 
\usepackage{listings}
\usepackage{xcolor}

\definecolor{codeblue}{rgb}{0,0,0.8}
\definecolor{codegreen}{rgb}{0,0.5,0}
\definecolor{codegray}{rgb}{0.4,0.4,0.4}
\definecolor{backcolour}{rgb}{0.97,0.97,0.97}

\lstdefinelanguage{RTA}{
    keywords={init, int, clock, inv, if, then, disabled, AND, OR, aut, name},
    sensitive=true,
    morecomment=[l]{//},
}

\lstset{
    language=RTA,
    backgroundcolor=\color{backcolour},
    basicstyle=\ttfamily\small,
    keywordstyle=\color{codeblue}\bfseries,
    commentstyle=\color{codegray}\itshape,
    stringstyle=\color{codegreen},
    numbers=left,
    numberstyle=\tiny\color{codegray},
    stepnumber=1,
    numbersep=8pt,
    showstringspaces=false,
    breaklines=true,
    frame=lines,
    xleftmargin=15pt,
    belowskip=0pt,
    aboveskip=10pt,
    literate={-->}{{\texttt{-->}}}3 {->>}{{\texttt{->>}}}3 {--!}{{\texttt{--!}}}3 {--\#--}{{\texttt{--\#--}}}5 {---->}{{\texttt{---->}}}5 {':=}{{\texttt{':=}}}3
}

\begin{document}

\begin{center}
    {\huge \textbf{Guia de Modelação e Verificação RTA}} \\
    \vspace{2mm}
    {\large Documentação Técnica -- \today}
\end{center}



\section{A Sintaxe da Linguagem RTA}

\subsection{Declarações Globais}
\begin{table}[h]
\centering\small
\renewcommand{\arraystretch}{1.2}
\begin{tabularx}{\textwidth}{l l X}
\toprule
\textbf{Sintaxe} & \textbf{Exemplo} & \textbf{Descrição} \\ \midrule
\texttt{init} & \texttt{init s0;} & Define o estado inicial do sistema. \\
\texttt{int} & \texttt{int c = 0;} & Declara uma variável inteira e o seu valor inicial. \\
\texttt{clock} & \texttt{clock t1;} & Declara relógios (iniciam sempre a 0). \\
\texttt{inv} & \texttt{inv s: t <= 5;} & Associa uma condição invariante a um estado. \\
\texttt{aut} & \texttt{\{ init s0; s0 --> s1: act \}} & Autómatos, permitem encapsular lógica. \\  \bottomrule
\end{tabularx}
\end{table}

\subsection{Arestas e Reconfiguração}
A estrutura base é: \texttt{origem <SETA> destino <ATRIBUTOS...>;}
\begin{table}[h]
\centering\small
\renewcommand{\arraystretch}{1.2}
\begin{tabularx}{\textwidth}{l X}
\toprule
\textbf{Seta} & \textbf{Comportamento Reativo} \\ \midrule
\texttt{-->} & Transição simples de estado. \\
\texttt{->>} & \textbf{Ativação:} O disparo ativa o rótulo alvo. \\
\texttt{--!} & \textbf{Desativação:} O disparo desativa o rótulo alvo. \\
\texttt{----> } & \textbf{Dependência:} Ativa o alvo e desativa-se a si próprio (atalho). \\ \bottomrule
\end{tabularx}
\end{table}

\begin{itemize}[noitemsep, topsep=0pt]
    \item \textbf{Rótulo:} \texttt{: nome} (ex: \texttt{s0 --> s1: clicar})
    \item \textbf{Guarda/Atualização:} \texttt{if (cond) then \{ var' := expr \}}
    \item \textbf{Estado Inicial:} \texttt{disabled} (começa inativa).
\end{itemize}


\section{Modelação na Prática}


\noindent
\begin{minipage}[t]{0.48\textwidth}
\begin{lstlisting}[caption={Modularidade e Escopo}]
name EX1
aut Motor {
  init Parado
  Parado --> Andamento: ligar
}
\end{lstlisting}
\end{minipage}
\hfill
\begin{minipage}[t]{0.48\textwidth}
\begin{lstlisting}[caption={Grafos Reativos}]
name EX2
init Off
Off --> On: ligar
On  --> Off: desligar
ligar --! ligar
desligar ->> ligar
\end{lstlisting}
\end{minipage}
\hfill
\begin{minipage}{0.48\textwidth}
\begin{lstlisting}[caption={Guardados}]
name EX3
int c = 0
init Ready
Ready --> Ready: step if (c < 3) then {
    c' := c + 1
}
\end{lstlisting}
\end{minipage}
\begin{minipage}{0.48\textwidth}
\begin{lstlisting}[caption={Temporizados}]
name EX4
clock t
init Idle
inv Work: t <= 10
Idle --> Work: start t' := 0
Work --> Idle: tout if (t >= 10)
\end{lstlisting}
\end{minipage}







\section{Verificação PDL (Lógica Proposicional Dinâmica)}

A plataforma traduz fórmulas visuais para lógica de backend. 

\subsection{Operadores Especiais}
\begin{itemize}[noitemsep]
    \item \textbf{Delay no PDL:} É possível verificar estados temporais futuros. \\
    Exemplo: \texttt{<delay:5>true} verifica se o sistema pode sobreviver 5 unidades de tempo.
    \item \textbf{Proposições de Estado:} \texttt{s1} é verdade se o sistema estiver no estado \texttt{s1}.
    \item \textbf{Proposições de Dados:} \texttt{c == 2} verifica o valor das variáveis.
\end{itemize}

\textbf{Exemplos Práticos:}
\begin{enumerate}[label=\arabic*., noitemsep]
    \item \textbf{Alcançabilidade:} $\langle \text{turn\_on} \rangle \text{S\_On}$
    \item \textbf{Lógica de Reset:} $[ \text{turn\_off} ] \langle \text{turn\_on} \rangle \text{true}$
    \item \textbf{Verificação de Variável:} $\langle \text{step; step} \rangle [\text{counter} == 2]$
    \item \textbf{Ausência de Bloqueio:} $\langle \text{move\_left} + \text{move\_right} \rangle \text{true}$
\end{enumerate}

\subsection{Exemplos de Verificação}
\begin{enumerate}[label=\arabic*., noitemsep]
    \item \textbf{Reconfiguração:} \texttt{<ligar> ligar\_active == 0} (Após ligar, a ação de ligar fica desativada?)
\end{enumerate}

\section{Exportação e Análise}
\begin{itemize}
    \item \textbf{Uppaal (TGRG):} Gera um modelo XML compatível com o Uppaal, mapeando a reconfiguração para uma estrutura de dados de matriz (\texttt{Edge A} e \texttt{Hyperedge L}).
    \item \textbf{Tradução GLTS:} Converte o modelo reconfigurável num modelo de estados padrão (plano), facilitando a leitura de ferramentas externas.
\end{itemize}





\end{document}